\section*{Introduction}

Psoriasis has a complex genetic architecture involving innate and adaptive immune pathways.\textsuperscript{\cite{tsoi2012psoriasis}} The present screen also included dermatitis and eczema, vitiligo, and alopecia areata outcomes to examine genetic tractability across distinct dermatologic phenotypes. Biologics and small molecules directed at the IL-17, IL-23, TNF, JAK, and TYK2 pathways have transformed dermatologic care, but response heterogeneity and failed development programs continue to motivate better target-validation strategies. The selective TYK2 inhibitor deucravacitinib is now an established treatment for moderate-to-severe plaque psoriasis.\textsuperscript{\cite{armstrong2023deucravacitinib}}

Drug-target Mendelian randomization (MR) uses germline variants near a gene to proxy perturbation of a therapeutic mechanism. Targets supported by human genetics are more likely to progress successfully through clinical development, although the validity of any individual MR estimate depends on instrument relevance, exclusion restriction, and appropriate handling of linkage disequilibrium (LD).\textsuperscript{\cite{nelson2015genetic,king2019drugtargets,schmidt2020drugtarget,burgess2023drugdiscovery}} Cis-expression quantitative trait loci (cis-eQTLs) provide one route to instrument gene expression, but regulatory effects are tissue-specific. GTEx therefore offers a useful framework for comparing skin and blood, tissues that capture different components of inflammatory skin disease biology.\textsuperscript{\cite{gtex2020,ongen2017causaltissues}}

A practical limitation is often overlooked: many clinically important genes do not have a sufficiently strong cis-eQTL in the tissue being studied. In that setting, a target is untestable rather than genetically unsupported. Conflating the absence of an instrument with evidence of no effect can distort target prioritization.\textsuperscript{\cite{gill2024pitfalls,sanderson2022mr}} This concern is especially relevant to cytokines, receptors, intracellular kinases, and barrier genes whose expression may be cell-state-specific or poorly captured in bulk tissue.

The TYK2 locus provides a useful positive-control benchmark. TYK2 coding variation, including P1104A, is associated with autoimmune protection, and pharmacologic TYK2 inhibition is effective in psoriasis.\textsuperscript{\cite{steri2025tyk2,armstrong2023deucravacitinib}} Previous MR and colocalization studies have also implicated TYK2 in psoriasis.\textsuperscript{\cite{guo2024tyk2,yuan2023tyk2}} Recovery of a TYK2 psoriasis signal can therefore test whether a tissue-aware pipeline detects an established locus. Importantly, however, the direction of an expression-based instrument cannot automatically be equated with the direction of kinase inhibition because a locus may contain coding and regulatory effects with different molecular consequences.

We therefore screened a curated panel of 19 clinically relevant dermatologic immune pathway genes across sun-exposed skin, non-sun-exposed skin, and whole blood. The panel deliberately included direct drug targets, target-family members, pathway genes, and disease-relevant controls rather than 19 direct targets of marketed drugs. We treated instrument availability as a primary result, performed a threshold-defined genome-wide-significant analysis and a separately corrected relaxed-threshold analysis, and evaluated TYK2 with locus inspection, external association data, and tissue-resolved colocalization. The study design is shown in Figure~1.
